\documentclass{zhvt-classic}

% 使用本地源流明体（GenRyuMin2 TW）按文件路径加载，不依赖系统字体名
% 版框（内框／界栏／书口侧栏／鱼尾）改用与版心印章同一印泥棕色：
%   印章图片 yshs.png 实色 = RGB(135,68,52) = HTML 874434
\usepackage{xcolor}
\usepackage{tikz}
% sealbrown 双注册：xcolor 供 TikZ/pgf，L3 供 cls 版框
\definecolor{sealbrown}{HTML}{874434}
\ExplSyntaxOn
\color_set:nnn { sealbrown } { HTML } { 874434 }
\ExplSyntaxOff

\zhvtset{
  grid color       = sealbrown,
  main font        = GenRyuMin2TW-SB,
  main font path   = C:/Users/Administrator/texmf/fonts/opentype/public/GenRyuMin2TW/,
jiazhu font      = GenRyuMin2TW-M,
  jiazhu font path = C:/Users/Administrator/texmf/fonts/opentype/public/GenRyuMin2TW/,
  % 页码框容纳 4 字：正常页「九十七」+前置空 1 格=4 行恰好；
  % 100+ 页「一百零一」共 5 行超出，末字「一/二/三」被框裁掉（2026-10-04 实测），故放宽到 5
  shukou page number chars = 5,
  % 书口书名框高从默认 6 字减到 4 字，把省出的 2 字高拨给 chapter 框，
  % 使 6 字书名+卷名放进 9 字高的 chapter 框后不再压线（2026-10-05）
  shukou title chars = 4,
}

\title{史记}[白起王翦列传]
\author{汉·司马迁}
\maker{闲叟漫记·重製}


% ===== 索隐述赞样式（2026-11-08 四次定稿 · 愚大人亲手改定参数）=====
% ⚠ 竖排前提：自定义字体家族必须带 RawFeature={vertical:+vert}；
%   simhei.ttf 加 +vhal 会整体加载失败静默回退，只写 +vert；simfang 同理。
\newCJKfontfamily{\shuzanfont}[Path=C:/Windows/Fonts/,RawFeature={vertical:+vert}]{simhei.ttf}
\newCJKfontfamily{\shuzanbody}[Path=C:/Windows/Fonts/,RawFeature={vertical:+vert}]{simfang.ttf}
% 题头：152x30pt 圆角 6pt 棕底白字；baseline (0,2) 使整框落入界行格子不压版框线；
% 白字节点 x 不补偿、y +8.2pt（像素级墨迹实测 dx/dy≈0，L/R 留白约 20pt）。
\newcommand{\shuzanhead}[1]{%
  \begin{tikzpicture}[baseline={(0,2)}]
    \fill[sealbrown, rounded corners=6pt]
      (-76pt,-15pt) rectangle (76pt,15pt);
    \node[text=white, inner sep=0pt, anchor=center] at (0pt,8.2pt)
      {\shuzanfont\fontsize{23}{23}\selectfont #1};
  \end{tikzpicture}%
}
% ★ 修正书口卷名位置：chapter 框起始位置改回 cls 默认值，避免全局影响。
%   只在每个卷首页通过 \juan 命令里的空命令微调该页位置（安全，不会跨页污染）。
\ExplSyntaxOn
\AtBeginDocument{
  \dim_set:Nn \l__zhvt_shukou_chapter_title_start_dim
    { \l__zhvt_shukou_title_height_dim + \l__zhvt_font_size_dim * 2 }
  \dim_set:Nn \l__zhvt_shukou_page_number_start_dim
    { \l__zhvt_shukou_chapter_title_start_dim + \l__zhvt_shukou_chapter_title_height_dim }
}
\ExplSyntaxOff

\begin{document}

% ★★ 版心（书口）排布覆写：把书名放到鱼尾之下（2026-10-04）
%
%   cls 默认版心自上而下：书名（title 静态框，在鱼尾之上）→ 鱼尾 → 篇名（chapter 动态框）→ 页码
%   本模板改为：          鱼尾（位置不动）→ 书名 → 篇名 → 页码分割线 → 页码
%
%   手法：把 cls 放在 title 静态框里的书名，改放进 chapter 动态框（鱼尾之下）；
%         title 静态框留空。书名取 \l__zhvt_title_str，即 \title{...} 的第一参数，
%         所以换书时只改 \title，不必在此重复写书名。
%
%   ⚠ 三条必须遵守的写法（2026-10-04 实测）：
%     ① 汉字必须用 \hbox{...} 包住。flowfram 会把框内容装进一个宽 14.4mm（= 一个 line_sep
%        = 书口栏宽）的纵向 minipage 并按该宽度折行，而整页靠 PDF /Rotate 90 显示、
%        框内文字本身不变向 —— 不包 \hbox 时「围炉夜话」4 字会被逐字折成 4 行，
%        在版面上就成了横排 4 字并横向溢出版框。
%     ② 竖排下**一个 \hbox 就是一整列**。书名与篇名要同处一列，必须放进同一个 \hbox；
%        分成两个 \hbox 时，后一个会被当成新的一列、甩到隔壁栏去
%        （实测「序」跑到书口栏外侧 14.4mm 的正文列上）。中间的 \quad 在框内是横向空白，
%        显示出来正好是书名与篇名之间的纵向间隔。
%     ③ ★ 内容盒要写成 **与框等宽的定宽盒** `\hbox to \l__zhvt_shukou_chapter_title_height_dim{…\hfil}`。
%        实测：chapter 框是把内容盒沿纵向**居中**摆放的，内容宽度一变，书名就跟着漂：
%        自然宽度的 \hbox（书名 38.4mm）落在 x=157.2mm，加篇名后（57.6mm）落在 x=152.4mm，
%        目录页（无篇名）与正文页的书名就差半格。做成等宽盒后，无论有没有篇名，
%        内容固定从框顶起排，实测两页均落在 x=142.8mm，不再漂。
\ExplSyntaxOn
% 把书口 chapter 框（书名+篇名）起始位置下移 1 字高，避免 4 字卷名顶到鱼尾
% ⚠ cls 的 \__zhvt_calculate_dimensions: 在 \AtBeginDocument 里执行，会重算该 dim；
%   必须把覆写挂到 begin document 之后，否则会被二次覆盖（2026-10-05 实测）。
% (旧方案已移到导言区，直接修改 chapter 框起始位置)
% 卷标题命令 #1=书口篇名(鱼尾) #2=正文卷名（含卷序，卷名内空格用 ~）
%
% ★★ 目录页两项修订（2026-10-05，愚大人指出）★★
%   ① 页码起点统一对齐：
%      原状 —— 目录项写成 \par\hspace{4em}卷名\hspace{4em}页码，页码起点 = 卷名实宽 + 4em。
%      卷名实宽随字数变：「卷一~梁惠王上」6 字 +~、「卷七~离娄上」5 字 +~，
%      差 1 个字位 → 卷七~卷十 4 条的页码整体比其余 10 条高 27.21pt（一个字），看着上下不齐。
%      改法 —— 卷名补全角空格（宽 = 1 字 = 27.215pt）补足到 **7 字符位**再写进目录，
%      14 条页码起点遂全部落在同一条竖线上（起点对齐；末字随页码字数自然参差）。
%      ⚠ 只对目录里的副本 \l_tmpc_tl 补，正文卷名仍用 \l_tmpa_tl 原样 —— 正文像素不受影响。
%   ② 卷名整体上移半字位：前导 \hspace 由 4em 改 3.5em。
%      实测原状目录块在版框内偏下（天头留白 139.6pt / 地脚留白 110.7pt），
%      上移 0.5em = 13.6pt 后上下基本均分，视觉居中。
\NewDocumentCommand { \juan } { mm } {
	\clearpage
	\tl_set:Nn \l__zhvt_chapter_title_tl { #1 }
	\tl_set:Nn \l_tmpa_tl { #2 }
	\tl_set_eq:NN \l_tmpc_tl \l_tmpa_tl
	\int_step_inline:nn { 7 - \str_count:n { #2 } }
	  { \tl_put_right:Nn \l_tmpc_tl { 　 } }
	\tl_set:Nx \l_tmpb_tl { \zhnum { page } }
	\tl_put_left:Nn \l_tmpb_tl { \hspace { 4em } }
	\tl_put_left:Nn \l_tmpb_tl { \l_tmpc_tl }
	\tl_put_left:Nn \l_tmpb_tl { \par \hspace { 3.5em } }
	\addtocontents { toc } { \l_tmpb_tl }
	\skip_horizontal:n { \l__zhvt_font_size_dim * 4 }
	\l_tmpa_tl
	\par
	\thispagestyle { empty }
}
\cs_set_protected:Npn \__zhvt_set_shukou_content: {
	% ★ 鱼尾上方留白（title 静态框）嵌入闲叟钤印：
	%   整页 /Rotate 90，故此处把图用 \rotatebox{90} 转正为竖章；
	%   尺寸卡在版心条幅内（书内约 7mm 宽 × 33.8mm 高，留余量不压鱼尾）。
	\setstaticcontents * { title } {
		\vbox_to_ht:nn { \l__zhvt_baseline_skip_dim } {
			\vskip 14.2pt
			\hbox_to_wd:nn { \l__zhvt_shukou_title_height_dim } {
				\hfil
				\rotatebox [ origin = c ] { 90 } { \includegraphics [ width = 7mm , height = 33.8mm ] { yshs.png } }
				\hfil
			}
		}
	}
% ★ 版心（书口）统一改法（2026-11-16 愚大人定）：18pt＋框内居中＋书口文字同版框色
%   ① 原 27.2pt 时 7 字长篇名（如「廉颇蔺相如列传」）末字溢出 chapter 框、压住页码；
%      改 \fontsize{18pt}{18pt}，前后 \hfil 令字串在框内居中（不随篇名字数漂移）。
%   ② \vbox{\vskip 8.2pt} 定量下移，抵消缩字号后栏宽方向 8.2pt 上偏（像素实测 y 中心 543.3 = 原版）。
%      ⚠ \raisebox / \vspace* 在 flowfram 框内无效；换字号须同步换补偿：16pt→10.1、19pt→7.4。
%   ③ chapter 与 page~number 均置 \color{sealbrown}（HTML 874434），与版框／界行同色；正文仍黑。
%   ⚠ 页码字位、钤印、鱼尾、界行、正文一律未动；五部合订本同此口径。
	\setdynamiccontents * { chapter }
	{ \color{sealbrown} \vbox { \vskip 8.2pt \hbox to \l__zhvt_shukou_chapter_title_height_dim
		{ \hfil \fontsize{18pt}{18pt}\selectfont \l__zhvt_title_str \quad \l__zhvt_chapter_title_tl \hfil } } }
	\setdynamiccontents * { page~number } { \color{sealbrown} \quad \zhnum { page } }
}
\__zhvt_set_shukou_content:
\ExplSyntaxOff

\maketitle

% 单篇，无多卷目录，不渲染目录页

\juan{白起王翦列传}{卷七十三~白起王翦列传第十三}
\hspace*{2em}白起者。郿人也。善用兵。事秦昭王。昭王十三年。而白起为左庶长。将而击韩之新城。是岁。穰侯相秦。举任鄙以为汉中守。其明年。白起为左更。攻韩。魏於伊阙。斩首二十四万。又虏其将公孙喜。拔五城。起迁为国尉。涉河取韩安邑以东。到乾河。明年。白起为大良造。攻魏。拔之。取城小大六十一。明年。起与客卿错攻垣城。拔之。後五年。白起攻赵。拔光狼城。後七年。白起攻楚。拔鄢。邓五城。其明年。攻楚。拔郢。烧夷陵。遂东至竟陵。楚王亡去郢。东走徙陈。秦以郢为南郡。白起迁为武安君。武安君因取楚。定巫。黔中郡。昭王三十四年。白起攻魏。拔华阳。走芒卯。而虏三晋将。斩首十三万。与赵将贾偃战。沈其卒二万人於河中。昭王四十三年。白起攻韩陉城。拔五城。斩首五万。四十四年。白起攻南阳太行道。绝之。\\
\hspace*{2em}四十五年。伐韩之野王。野王降秦。上党道绝。其守冯亭与民谋曰。﹁郑道已绝。韩必不可得为民。秦兵日进。韩不能应。不如以上党归赵。赵若受我。秦怒。必攻赵。赵被兵。必亲韩。韩赵为一。则可以当秦。﹂因使人报赵。赵孝成王与平阳君。平原君计之。平阳君曰。﹁不如勿受。受之。祸大於所得。﹂平原君曰。﹁无故得一郡。受之便。﹂赵受之。因封冯亭为华阳君。\\
\hspace*{2em}四十六年。秦攻韩缑氏。蔺。拔之。\\
\hspace*{2em}四十七年。秦使左庶长王龁攻韩。取上党。上党民走赵。赵军长平。以按据上党民。四月。龁因攻赵。赵使廉颇将。赵军士卒犯秦斥兵。秦斥兵斩赵裨将茄。六月。陷赵军。取二鄣四尉。七月。赵军筑垒壁而守之。秦又攻其垒。取二尉。败其阵。夺西垒壁。廉颇坚壁以待秦。秦数挑战。赵兵不出。赵王数以为让。而秦相应侯又使人行千金於赵为反间。曰。﹁秦之所恶。独畏马服子赵括将耳。廉颇易与。且降矣。﹂赵王既怒廉颇军多失亡。军数败。又反坚壁不敢战。而又闻秦反间之言。因使赵括代廉颇将以击秦。秦闻马服子将。乃阴使武安君白起为上将军。而王龁为尉裨将。令军中有敢泄武安君将者斩。赵括至。则出兵击秦军。秦军详败而走。张二奇兵以劫之。赵军逐胜。追造秦壁。壁坚拒不得入。而秦奇兵二万五千人绝赵军後。又一军五千骑绝赵壁间。赵军分而为二。粮道绝。而秦出轻兵击之。赵战不利。因筑壁坚守。以待救至。秦王闻赵食道绝。王自之河内。赐民爵各一级。发年十五以上悉诣长平。遮绝赵救及粮食。\\
\hspace*{2em}至九月。赵卒不得食四十六日。皆内阴相杀食。来攻秦垒。欲出。为四队。四五复之。不能出。其将军赵括出锐卒自搏战。秦军射杀赵括。括军败。卒四十万人降武安君。武安君计曰。﹁前秦已拔上党。上党民不乐为秦而归赵。赵卒反覆。非尽杀之。恐为乱。﹂乃挟诈而尽阬杀之。遗其小者二百四十人归赵。前後斩首虏四十五万人。赵人大震。\\
\hspace*{2em}四十八年十月。秦复定上党郡。秦分军为二。王龁攻皮牢。拔之。司马梗定太原。韩。赵恐。使苏代厚币说秦相应侯曰。﹁武安君禽马服子乎。﹂曰。﹁然。﹂又曰。﹁即围邯郸乎。﹂曰。﹁然。﹂﹁赵亡则秦王王矣。武安君为三公。武安君所为秦战胜攻取者七十馀城。南定鄢。郢。汉中。北禽赵括之军。虽周。召。吕望之功不益於此矣。今赵亡。秦王王。则武安君必为三公。君能为之下乎。虽无欲为之下。固不得已矣。秦尝攻韩。围邢丘。困上党。上党之民皆反为赵。天下不乐为秦民之日久矣。今亡赵。北地入燕。东地入齐。南地入韩。魏。则君之所得民亡几何人。故不如因而割之。无以为武安君功也。﹂於是应侯言於秦王曰。﹁秦兵劳。请许韩。赵之割地以和。且休士卒。﹂王听之。割韩垣雍。赵六城以和。正月。皆罢兵。武安君闻之。由是与应侯有隙。\\
\hspace*{2em}其九月。秦复发兵。使五大夫王陵攻赵邯郸。是时武安君病。不任行。四十九年正月。陵攻邯郸。少利。秦益发兵佐陵。陵兵亡五校。武安君病愈。秦王欲使武安君代陵将。武安君言曰。﹁邯郸实未易攻也。且诸侯救日至。彼诸侯怨秦之日久矣。今秦虽破长平军。而秦卒死者过半。国内空。远绝河山而争人国都。赵应其内。诸侯攻其外。破秦军必矣。不可。﹂秦王自命。不行。乃使应侯请之。武安君终辞不肯行。遂称病。\\
\hspace*{2em}秦王使王龁代陵将。八九月围邯郸。不能拔。楚使春申君及魏公子将兵数十万攻秦军。秦军多失亡。武安君言曰。﹁秦不听臣计。今如何矣。﹂秦王闻之。怒。彊起武安君。武安君遂称病笃。应侯请之。不起。於是免武安君为士伍。迁之阴密。武安君病。未能行。居三月。诸侯攻秦军急。秦军数卻。使者日至。秦王乃使人遣白起。不得留咸阳中。武安君既行。出咸阳西门十里。至杜邮。秦昭王与应侯群臣议曰。﹁白起之迁。其意尚怏怏不服。有馀言。﹂秦王乃使使者赐之剑。自裁。武安君引剑将自刭。曰。﹁我何罪于天而至此哉。﹂良久。曰。﹁我固当死。长平之战。赵卒降者数十万人。我诈而尽阬之。是足以死。﹂遂自杀。武安君之死也。以秦昭王五十年十一月。死而非其罪。秦人怜之。乡邑皆祭祀焉。\\
\hspace*{2em}王翦者。频阳东乡人也。少而好兵。事秦始皇。始皇十一年。翦将攻赵阏与。破之。拔九城。十八年。翦将攻赵。岁馀。遂拔赵。赵王降。尽定赵地为郡。明年。燕使荆轲为贼於秦。秦王使王翦攻燕。燕王喜走辽东。翦遂定燕蓟而还。秦使翦子王贲击荆。荆兵败。还击魏。魏王降。遂定魏地。\\
\hspace*{2em}秦始皇既灭三晋。走燕王。而数破荆师。秦将李信者。年少壮勇。尝以兵数千逐燕太子丹至於衍水中。卒破得丹。始皇以为贤勇。於是始皇问李信。﹁吾欲攻取荆。於将军度用几何人而足。﹂李信曰。﹁不过用二十万人。﹂始皇问王翦。王翦曰。﹁非六十万人不可。﹂始皇曰。﹁王将军老矣。何怯也。李将军果势壮勇。其言是也。﹂遂使李信及蒙恬将二十万南伐荆。王翦言不用。因谢病。归老於频阳。李信攻平与。蒙恬攻寝。大破荆军。信又攻鄢郢。破之。於是引兵而西。与蒙恬会城父。荆人因随之。三日三夜不顿舍。大破李信军。入两壁。杀七都尉。秦军走。\\
\hspace*{2em}始皇闻之。大怒。自驰如频阳。见谢王翦曰。﹁寡人以不用将军计。李信果辱秦军。今闻荆兵日进而西。将军虽病。独忍弃寡人乎。﹂王翦谢曰。﹁老臣罢病悖乱。唯大王更择贤将。﹂始皇谢曰。﹁已矣。将军勿复言。﹂王翦曰。﹁大王必不得已用臣。非六十万人不可。﹂始皇曰。﹁为听将军计耳。﹂於是王翦将兵六十万人。始皇自送至灞上。王翦行。请美田宅园池甚众。始皇曰。﹁将军行矣。何忧贫乎。﹂王翦曰。﹁为大王将。有功终不得封侯。故及大王之乡臣。臣亦及时以请园池为子孙业耳。﹂始皇大笑。王翦既至关。使使还请善田者五辈。或曰。﹁将军之乞贷。亦已甚矣。﹂王翦曰。﹁不然。夫秦王怚而不信人。今空秦国甲士而专委於我。我不多请田宅为子孙业以自坚。顾令秦王坐而疑我邪。﹂\\
\hspace*{2em}王翦果代李信击荆。荆闻王翦益军而来。乃悉国中兵以拒秦。王翦至。坚壁而守之。不肯战。荆兵数出挑战。终不出。王翦日休士洗沐。而善饮食抚循之。亲与士卒同食。久之。王翦使人问军中戏乎。对曰。﹁方投石超距。﹂於是王翦曰。﹁士卒可用矣。﹂荆数挑战而秦不出。乃引而东。翦因举兵追之。令壮士击。大破荆军。至蕲南。杀其将军项燕。荆兵遂败走。秦因乘胜略定荆地城邑。岁馀。虏荆王负刍。竟平荆地为郡县。因南征百越之君。而王翦子王贲。与李信破定燕。齐地。\\
\hspace*{2em}秦始皇二十六年。尽并天下。王氏。蒙氏功为多。名施於後世。\\
\hspace*{2em}秦二世之时。王翦及其子贲皆已死。而又灭蒙氏。陈胜之反秦。秦使王翦之孙王离击赵。围赵王及张耳钜鹿城。或曰。﹁王离。秦之名将也。今将彊秦之兵。攻新造之赵。举之必矣。﹂客曰。﹁不然。夫为将三世者必败。必败者何也。必其所杀伐多矣。其後受其不祥。今王离已三世将矣。﹂居无何。项羽救赵。击秦军。果虏王离。王离军遂降诸侯。\\
\hspace*{2em}太史公曰。鄙语云﹁尺有所短。寸有所长﹂。白起料敌合变。出奇无穷。声震天下。然不能救患於应侯。王翦为秦将。夷六国。当是时。翦为宿将。始皇师之。然不能辅秦建德。固其根本。偷合取容。以至筊身。及孙王离为项羽所虏。不亦宜乎。彼各有所短也。\\

% ===== 索隐述赞（司马贞《史记索隐》）· 统一样式 =====
% 题头「索隐述赞曰」五字并入反白框；正文每四句一列，定宽盒 \hfil 居列中；
% \fontsize{24}{41} 列距对齐全书界行（121+41k）；\clearpage 使述赞独占一页。
\clearpage
\shuzanhead{索隐述赞曰}\linebreak
{\shuzanbody\fontsize{24}{41}\selectfont
\hbox to \hsize{\hfil 白起。王翦。俱善用兵。递为秦将。\hfil}\linebreak
\hbox to \hsize{\hfil 拔齐破荆。赵任马服。长平遂阬。楚陷李信。\hfil}\linebreak
\hbox to \hsize{\hfil 霸上卒行。贲。离继出。三代无名。\hfil}%
\par}
\end{document}
